\section{Proof of the limiting oracle characterization}
\label{sec:oracle-limit-proofs}

We prove \cref{thm:oracle-limits} uniformly over the original rate-and-stop
domain. Throughout, $s,r,\nu\geq0$ are fixed and finite,
$b=q(1-q)$, $w=(1-q)^2$, $v=s/4+r$, and $T=M_qb\in[T_0,T_1]$ for fixed
$0<T_0<T_1<\infty$. Constants in the $O(\cdot)$ terms may depend on these
fixed quantities, but not on the rate, stop, or cap in the stated ranges.

\subsection{Exact moments and large-rate bounds}

Rescale the common moment in \cref{eq:predictive-moments} by setting
$G_k=C_k/4$ and restore the arbitrary noise variance $\nu$. The recurrence is
\begin{equation}
  \begin{aligned}
    G_{k+1}&=(1-u)^2G_k+u^2bH_k+u^2\nu,\\
    H_{k+1}&=D_b(\alpha)H_k+b\alpha^2(G_k+\nu),\\
    D_b(\alpha)&=1-b\alpha(2-\alpha+3b\alpha),
  \end{aligned}
  \label{eq:limit-moments}
\end{equation}
with $u=b\alpha$ for S and $u=\alpha/4$ for U. Initially
$(G_0,H_0)=(s/4,r)$, and $R_{R,k}=(G_k+wH_k)/2$. All coefficients are
nonnegative, since
$D_b=(1-b\alpha)^2+b\alpha^2(1-4b)\geq0$.

For S, let $W_k=G_k+H_k$. Summing the two equations gives
\[
  W_{k+1}\geq\lambda_b(\alpha)W_k+b\alpha^2\nu,
  \qquad
  \lambda_b(\alpha)=1-2b\alpha+b\alpha^2(1-3b+b^2).
\]
The coefficient of $H_k$ is exactly $\lambda_b$, and that of $G_k$ exceeds it
by $b^2\alpha^2(4-b)$. The exact noise term is $b\alpha^2(1+b)\nu$.
Also $\lambda_b\geq1-b/(1-3b+b^2)>0$ for $0<b\leq1/4$. Consequently,
\[
  R^S_{R,k}(\alpha)\geq\frac w2\left[
    \lambda_b^kv+b\alpha^2\nu\sum_{j=0}^{k-1}\lambda_b^j\right].
\]
In particular, for $\alpha\geq c_S(b):=2/(1-3b+b^2)$, this risk is at least
$wv/2$ at every stop, including initialization.

For U, positivity gives the analogous bound
\[
  R^U_{R,k}(\alpha)\geq\frac w2\left[
    D_b^kr+b\alpha^2\nu\sum_{j=0}^{k-1}D_b^j\right].
\]
If $\alpha\geq c_U(b):=2/(1-3b)$, then $D_b\geq1$ and the risk is at least
$wr/2$. Empty sums cover $k=0$. These bounds require no stability assumption
and hold even when rates diverge or stops are chosen to offset their growth.
For sufficiently small $q$, both cutoffs are below three and $1-w\leq4b$.

\subsection{Uniform approximation for standardization}

On $\alpha\in[0,3]$, the affine transition for $(G,H,1)$ under S is
$I+b\mathcal L_\alpha+O(b^2)$ uniformly, where
\[
  \mathcal L_\alpha=
  \begin{pmatrix}
    -2\alpha&0&0\\
    \alpha^2&-d(\alpha)&\alpha^2\nu\\
    0&0&0
  \end{pmatrix}.
\]
Its difference from $e^{b\mathcal L_\alpha}$ is $O(b^2)$, and both operators
have norm at most $1+O(b)$. Telescoping their powers over $k\leq T_1/b$
steps therefore gives an $O(b)$ error uniformly in rate and stop. At time
$t=kb$, the limiting moments solve
\[
  \dot G=-2\alpha G,\qquad
  \dot H=-d(\alpha)H+\alpha^2(G+\nu).
\]
Their sum solves
\[
  \dot W=-d(\alpha)W+\alpha^2\nu,\qquad W(0)=v.
\]
Define $\Psi_S(\alpha,t)=W(t)/2$. The risk weights differ from their limits
by $O(b)$, so
\begin{equation}
  R^S_{R,k}(\alpha)=\Psi_S(\alpha,kb)+O(b)
  \label{eq:limit-s-uniform}
\end{equation}
uniformly over the retained rates and all allowed stops. For $0<\alpha<2$,
solving the scalar equation gives $F_S$ with budget $t$. The flow also defines
the endpoints smoothly: $\Psi_S(0,t)=v/2$ and
$\Psi_S(2,t)=(v+4\nu t)/2$. A vanishing SD rate retains both initial
moments uniformly over all stops.

\subsection{Uniform approximation with the unscaled transient}

For U, retain $0\leq\alpha\leq c_U(b)$ and write
\[
  A=(1-\alpha/4)^2,\qquad
  h=\frac{\alpha}{8-\alpha},\qquad g=\nu h,\qquad c=bh,
  \qquad
  B=\frac{r+T\alpha^2(G_0+g+\nu)}{1-T\alpha^2c}.
\]
Here $0\leq D_b\leq1$, $h\leq3/5$, and
$T\alpha^2c\leq27T_1b/5$. The denominator is at least $1/2$ for small $b$.
With $H_{\max}=\max_{j\leq M_q}H_j$, summing the common recurrence gives
\[
  0\leq G_j-A^jG_0-(1-A^j)g\leq cH_{\max}.
\]
Using $G_j\leq G_0+g+cH_{\max}$ in the rare recurrence then gives
\[
  H_{\max}\leq r+T\alpha^2(G_0+g+\nu+cH_{\max}).
\]
Absorbing the last term yields
\[
  H_{\max}\leq B\leq2\{r+9T_1(G_0+8\nu/5)\},\qquad
  0\leq G_j-A^jG_0-(1-A^j)g\leq cB.
\]

Let $\widetilde H_0=r$ and
$\widetilde H_{j+1}=D_b\widetilde H_j+b\alpha^2(\nu+g)$. Since $D_b\leq1$,
\begin{align*}
  |H_k-\widetilde H_k|
  &\leq b\alpha^2\sum_{j<k}|G_j-g|\\
  &\leq b\{16h|G_0-g|+T\alpha^2hB\}=O(b).
\end{align*}
For $\alpha>0$, the geometric sum uses
$\alpha^2/(1-A)=16h$. At $\alpha=0$, both sides are zero. This cancellation
removes the apparent singularity in the common-coordinate relaxation time.

Let $H^\circ_\alpha$ solve
\[
  \dot H^\circ=-d(\alpha)H^\circ+\alpha^2(\nu+g),
  \qquad H^\circ(0)=r.
\]
The scalar affine map for $\widetilde H$ differs from this flow over time $b$
by $O(b^2)$ uniformly on the retained rate interval. The flow is uniformly
bounded for $t\leq T_1$, and errors propagate with multiplier $D_b\leq1$.
Telescoping gives $\widetilde H_k=H^\circ_\alpha(kb)+O(b)$. Hence
\begin{equation}
  R^U_{R,k}(\alpha)
    =\frac12\{A^kG_0+(1-A^k)g+H^\circ_\alpha(kb)\}+O(b)
    \label{eq:limit-u-uniform}
\end{equation}
uniformly over every retained rate and stop. Replacing the common term by $g$
uniformly would be false, including at initialization.

For clarity, this approximation also describes the boundary sequences. If
$\alpha_q\to0$, positivity and Bernoulli's inequality give
\[
  r-2rT\alpha_q\leq H_j\leq B=r+O(\alpha_q^2),\qquad
  0\leq G_j-A^jG_0\leq g+cB=O(\alpha_q).
\]
If $k_q\alpha_q\to\ell\in[0,\infty]$, the limiting risk is therefore
\begin{equation}
  R^U_{R,k_q}(\alpha_q)\longrightarrow
    \frac{s}{8}e^{-\ell/2}+\frac r2.
  \label{eq:limit-u-transient}
\end{equation}
For finite $\ell$, use $\log A=-\alpha_q/2+O(\alpha_q^2)$ and
$k_q\alpha_q^2\to0$. For infinite $\ell$, use
$A^{k_q}\leq e^{-k_q\alpha_q/2}$. Thus $b\ll\alpha_q\ll1$ at the full
cap learns the common signal and attains $r/2$.

If instead $\alpha_q\to\alpha_0\in(0,2]$ and $k_qb\to0$, the rare moment
tends to $r$. Along subsequences with $k_q=K$ fixed or $k_q\to\infty$, the
common moment tends, respectively, to
$A(\alpha_0)^KG_0+\{1-A(\alpha_0)^K\}g(\alpha_0)$ or $g(\alpha_0)$.
Both give risk at least $r/2$. These cases explain why substituting time zero
into the positive-time limit alone would miss some attainable transients.

\subsection{Reduction of the limiting trajectories}

Set $\Psi_U(\alpha,t)=\{g(\alpha)+H^\circ_\alpha(t)\}/2$. Both $\Psi_A$
are smooth through $\alpha=0,2$ when defined by their differential equations,
and $\Psi_U(0,t)=r/2$. For $0<\alpha<2$, each trajectory is monotone in
time. An increasing S trajectory stays above $v/2$, and an increasing U
trajectory stays above $(g+r)/2\geq r/2$. The respective baselines are the
values at rate zero in the limiting problems. A decreasing trajectory reaches
its minimum at the full cap $T$.

Rates $\alpha\in(1,2)$ are dominated by $2-\alpha\in(0,1)$ at the same time:
the decay $d(\alpha)$ is unchanged, while each noise term
\[
  g(\alpha),\qquad \frac{\nu\alpha}{2-\alpha},\qquad
  \frac{\alpha\{\nu+g(\alpha)\}}{2-\alpha}
\]
is nondecreasing in $\alpha$. At rate two the risks are
$(v+4\nu t)/2$ and $\{g(2)+r+4(\nu+g(2))t\}/2$, also above the baselines.
Consequently,
\begin{equation}
  \Psi_A(\alpha,t)\geq L_A(T),
  \qquad 0\leq\alpha\leq2,\quad 0\leq t\leq T.
  \label{eq:limit-trajectory-reduction}
\end{equation}
On $[0,3]\times[0,T_1]$, both flows and their rate derivatives are bounded,
as follows by differentiating their finite-dimensional matrix exponentials.
Since $c_A(b)-2=O(b)$, replacing a retained rate above two by two changes its
limiting risk by only $O(b)$.

\subsection{Lower bounds and recovery over the full oracle domain}

Combining \cref{eq:limit-s-uniform,eq:limit-trajectory-reduction} gives
$R^S_{R,k}(\alpha)\geq L_S(T)-O(b)$ at every retained rate and stop. Above
$c_S(b)$, the exact bound gives
$R^S_{R,k}\geq wv/2\geq L_S(T)-2bv$. The same lower bound therefore holds
over the entire S oracle domain.

For U, split the retained controls at $J_b=2\log(1/b)$. If
$k\alpha\geq J_b$, then $A^k\leq e^{-k\alpha/2}\leq b$, so
\cref{eq:limit-u-uniform,eq:limit-trajectory-reduction} give
$R^U_{R,k}(\alpha)\geq L_U(T)-O(b)$. If $k\alpha<J_b$, take $b$ small
enough that $1-2b\alpha\geq0$. Positivity gives
\[
  H_k\geq D_b^kr\geq(1-2b\alpha)^kr\geq r(1-2bk\alpha).
\]
It follows that
\[
  R^U_{R,k}(\alpha)\geq\frac r2-2br-rbJ_b
    \geq L_U(T)-2br-2rb\log(1/b).
\]
Above $c_U(b)$, the bound $R^U_{R,k}\geq wr/2\geq L_U(T)-2br$ suffices.
Thus the U oracle is bounded below by $L_U(T)-O(b\log(1/b))$. These bounds
permit taking infima directly, without requiring finite-prevalence minimizers
to exist or their controls to converge.

For matching upper bounds, let $\alpha_A(T)\in[0,1]$ minimize $F_A$.
Continuity ensures existence. For S, use $\alpha_S(T)$ at $k=M_q$;
\cref{eq:limit-s-uniform} gives risk $L_S(T)+O(b)$. For U, use the full cap
and the rate
\[
  \alpha_q=\max\left\{\alpha_U(T),\frac{2b\log(1/b)}{T}\right\}.
\]
This lies in $[0,1]$ for small $b$ and ensures $M_q\alpha_q\geq J_b$.
The common transient is $O(b)$, and $F_U$ has a bounded rate derivative on
$[0,1]\times[T_0,T_1]$. Hence
\[
  R^U_{R,M_q}(\alpha_q)=F_U(\alpha_q;T)+O(b)
    \leq L_U(T)+O(b\log(1/b)).
\]
This construction includes zero or vanishing minimizing rates and zero signal
or noise. Together with the lower bounds, it proves
\cref{eq:oracle-limit-errors}.

Finally, bounding the time derivatives of the objectives uniformly in rate and
taking minima yields
\[
  |L_S(T)-L_S(\tau)|\leq\frac{v+\nu}{2}|T-\tau|,\qquad
  |L_U(T)-L_U(\tau)|\leq\frac{r+8\nu/7}{2}|T-\tau|.
\]
Since $T\to\tau\in(0,\infty)$, both oracle limits follow. Rounding $\tau/b$
to an integer introduces an additional budget discrepancy of only $O(b)$.

\paragraph{Why the positive-budget assumption matters.}

The theorem does not assert uniformity for diverging signals, noise, or budgets.
Nor can $\tau>0$ be dropped without specifying the cap sequence. For
$s=4,r=0,\nu=1$, exact one-step minimization gives
\[
  R_U^*(1)=\frac12-\frac{1}{4(1+16bw)}\longrightarrow\frac14.
\]
If instead $M_q\to\infty$ and $M_qb\to0$, taking
$\alpha_q=M_q^{-1/2}$ gives risk tending to zero by the vanishing-rate bounds.
Both sequences have zero limiting rescaled budget, but different oracle limits.

\section{Variational decisions from the limiting risks}
\label{sec:oracle-variational}

\subsection{Unique optimal rates}

For $\nu>0$, write $x=v/\nu$, $\rho=r/\nu$, and, for $\alpha\in[0,1]$,
\[
  p=\frac{\alpha}{2-\alpha},\qquad
  h=\frac{\alpha}{8-\alpha},\qquad P=p(1+h),\qquad z=e^{-\tau d(\alpha)}.
\]
Then $L_S=\nu J_S/2$ and $L_U=\nu J_U/2$, where
\begin{equation}
  J_S(x,\tau)=\min_{[0,1]}\{xz+p(1-z)\},\qquad
  J_U(\rho,\tau)=\min_{[0,1]}\{\rho z+h+P(1-z)\}.
  \label{eq:dimensionless-oracles}
\end{equation}
Only the two signal-to-noise ratios and the rescaled budget determine the
preferred scaling. Multiplying $s,r,\nu$ by the same positive number scales
both risks by that number and leaves the optimal rates unchanged.

On $0<\alpha<1$, define
\[
  Q_S=p+\frac{p'(e^{\tau d}-1)}{\tau d'},\qquad
  Q_U=P+\frac{h'e^{\tau d}+P'(e^{\tau d}-1)}{\tau d'},
  \qquad d'=2(1-\alpha).
\]
Primes denote rate derivatives. Both functions increase strictly: $p,h,P$ have
positive first and second derivatives, and $e^{\tau d}$ and $1/d'$ increase
on this interval. Their endpoint limits are
\[
  Q_S(0)=0,\qquad Q_U(0)=\frac{1}{16\tau},\qquad
  Q_S(1-)=Q_U(1-)=\infty.
\]
The derivatives of the respective objectives in
\cref{eq:dimensionless-oracles} are $\tau d'z(Q_S-x)$ and
$\tau d'z(Q_U-\rho)$. Thus, if $v,\nu>0$, S has one minimizer in $(0,1)$,
given by $Q_S=x$. U has one minimizer in $(0,1)$ when
$r\tau>\nu/16$, given by $Q_U=\rho$, and otherwise its unique minimizer is
zero, including at equality. If $v=0$, S also minimizes uniquely at zero for
$\nu>0$. If $\nu=0$, choosing rate one gives
$L_S=ve^{-\tau}/2$ and $L_U=re^{-\tau}/2$; a zero signal makes the
corresponding objective identically zero.

The numerical solver in the source supplement brackets
these roots and multiplies the derivative by $z$ to avoid growing exponentials.
It evaluates the proved characterization in floating-point arithmetic; its
answers are not interval certificates.

\subsection{Allocation boundaries}

\Cref{cor:oracle-allocation} holds the total initial signal contribution $v$
fixed. Its boundary can be written
\begin{equation}
  \theta_*=
  \frac1x\max_{0\leq\alpha\leq1}
    \left[\{J_S(x,\tau)-h-P(1-z)\}e^{\tau d}\right].
  \label{eq:allocation-boundary}
\end{equation}
For each rate, the bracketed expression is the rare signal-to-noise ratio at
which that U trajectory equals $J_S$. U wins if at least one trajectory lies
below $J_S$, which explains the maximum. The lower and upper slope bounds used
in the corollary also give
\[
  \frac{e^{-\tau}}2(r_2-r_1)
  \leq L_U(r_2)-L_U(r_1)\leq\frac12(r_2-r_1).
\]

An equivalent comparison fixes $r,\nu,\tau>0$ and varies $s$. U's limiting
risk is constant, while
\[
  \frac{e^{-\tau}}8(s_2-s_1)
  \leq L_S(s_2)-L_S(s_1)\leq\frac18(s_2-s_1).
\]
At $s=0$, S wins by the same endpoint argument as in
\cref{cor:oracle-allocation}; its risk then grows without bound with $s$.
There is therefore exactly one $s_*>0$ giving equality, with S preferred below
it and U above it. Its variational expression and a bracket are
\[
  s_*=4\nu\max_{0\leq\alpha\leq1}
    \left[\{J_U(\rho,\tau)-p\}e^{\tau d}+p-\rho\right],
  \qquad 0<s_*\leq4r(e^\tau-1).
\]
The bracket follows from $L_S\geq ve^{-\tau}/2$ and $L_U\leq r/2$.
Either boundary can also be computed by scalar bisection of the optimized risk
difference. Uniqueness does not ensure good conditioning when $\tau$ is large.
These results concern the stated allocation paths and do not assert a unique
transition under every constraint on signal strength.

\subsection{Noise and budget regimes}
\label{sec:oracle-noise-budget}

For fixed signals and budget, both limiting risks are nondecreasing and concave
in $\nu$, because they are minima of affine functions with nonnegative noise
coefficients. Their difference need not be monotone. If $s>0$, zero noise
gives $L_S-L_U=se^{-\tau}/8>0$, so sufficiently small positive noise also
favors U. Once $\nu\geq16r\tau$, U has risk $r/2$. As $\nu\to\infty$,
the S minimizing rate must tend to zero, since rates bounded away from zero
have unbounded noise cost. Its signal term tends to $v/2$, while initialization
provides the matching upper bound. Hence
\[
  L_S\longrightarrow v/2,\qquad L_U\longrightarrow r/2,
  \qquad L_S-L_U\longrightarrow s/8.
\]
U also wins at sufficiently large noise when $s>0$. For positive common
signals in the rare-dominant family of \cref{thm:conference-predictive}, S wins
at noise variance $\nu=r$ and $\tau\in[1,4]$. Continuity then implies at least
two noise crossings. It does not imply exactly two, or crossings for every
signal pair. As a numerical illustration, $s=0.1,r=1,\tau=2$ gives:
\begin{center}
  \begin{tabular}{rrrr}
    \toprule
    $\nu$ & $L_U$ & $L_S$ & Preferred scaling\\
    \midrule
    0 & 0.067668 & 0.069359 & U\\
    0.1 & 0.109877 & 0.103572 & S\\
    1 & 0.256267 & 0.232599 & S\\
    100 & 0.500000 & 0.502306 & U\\
    \bottomrule
  \end{tabular}
\end{center}

Each limiting oracle risk is nonincreasing in $\tau$: the trajectory reduction
allows the earlier stopping choice under every larger cap. For fixed
$v,\nu>0$ and $\tau\downarrow0$, expanding uniformly on $\alpha\in[0,1]$
gives
\[
  L_S=\frac v2-\frac{\tau v^2}{2(v+\nu)}+O(\tau^2),\qquad
  \alpha_S\longrightarrow\frac{v}{v+\nu}.
\]
The first-order improvement in twice the risk is
$\tau\{2\alpha v-\alpha^2(v+\nu)\}$, a strictly concave quadratic.
For sufficiently small $\tau$, U has risk $r/2$ by its endpoint threshold.
Thus small positive rescaled budgets favor U when $s>0$. This statement takes
the prevalence limit first; the common coordinate has already been learned.

For fixed $r,\nu>0$ and finite $s\geq0$, the large-budget expansions are
\begin{equation}
  \begin{aligned}
    L_S&=\frac{\nu}{8\tau}
      \left[1+\log\!\left(\frac{4\tau v}{\nu}\right)\right]
      +o(\tau^{-1}),\\
    L_U&=\frac{5\nu}{32\tau}
      \left[1+\log\!\left(\frac{16\tau r}{5\nu}\right)\right]
      +o(\tau^{-1}).
  \end{aligned}
  \label{eq:large-budget-oracles}
\end{equation}
To justify them, the candidate $\alpha=(\log\tau)/(2\tau)$ gives risk
$O((\log\tau)/\tau)$. At any minimizer, the positive signal term forces
$z=O((\log\tau)/\tau)$. Under S, the noise term $\nu p(1-z)$ then bounds
$\alpha$ by the same order; under U, $\nu h$ does so directly. Thus
$\tau\alpha^2\to0$. Uniformly in this minimizing range, twice the risks are
\[
  ve^{-2\tau\alpha}+\frac{\nu\alpha}{2}+o(\tau^{-1}),\qquad
  re^{-2\tau\alpha}+\frac{5\nu\alpha}{8}+o(\tau^{-1}).
\]
The remainders are $O((\log\tau)^3/\tau^2)$ at both the exact minimizers and
the minimizers of these approximations. Minimizing
$Be^{-2\tau\alpha}+c\alpha$ gives
$\alpha=\log(2\tau B/c)/(2\tau)$ for large $\tau$, proving
\cref{eq:large-budget-oracles}. U pays $\nu\alpha/8$ in common-coordinate
noise in addition to $\nu\alpha/2$ in rare-coordinate noise, while S has only
the latter at leading order. Consequently $L_S/L_U\to4/5$.

For $s>0$, the small- and large-budget signs imply at least one budget crossing,
without asserting uniqueness. At $s=r=\nu=1$, numerical risks favor U at
$\tau=1$ and S at $\tau=10$. Convergence can be slow: for $s=10,r=0.1,\nu=2$,
U still wins at $\tau=10^5$, while S wins at $\tau=10^7$. The assumptions
matter: with $\nu=0$, U is always weakly better; with $r=0$, its limiting risk
is zero; with $s=0,r,\nu>0$, S wins at every finite positive budget.

The deterministic checks in the source supplement exercise the moment
approximations, bounds on unrestricted infima, scalar roots, boundaries, and
expansions. They supplement the proofs and are separate from
the finite-prevalence certificates.
